
\documentclass[a4paper,11pt]{report}

%\usepackage[landscape]{geometry}
\usepackage[cm-default]{fontspec}
\setromanfont{FreeSerif}
\setsansfont{FreeSans}
\setmonofont{FreeMono}
\usepackage[utf8x]{inputenc}
\usepackage{fontspec}
\usepackage{xunicode}
\usepackage{xltxtra}
\usepackage{xgreek}
\usepackage{amsmath}
\usepackage{unicode-math}
\usepackage{ulem}
\usepackage{color}
\usepackage{verbatim}
\usepackage{nopageno}
\usepackage{graphicx}
\usepackage{textpos}
\setlength{\TPHorizModule}{1cm}
\setlength{\TPVertModule}{1cm}
%\usepackage[colorgrid,texcoord]{eso-pic}
\usepackage[outline]{contour}
\usepackage{wrapfig}
\usepackage{url}
\usepackage{color}
\usepackage{tikz}
\usepackage{fancybox,fancyhdr}
\usepackage{subfigure}
\usepackage{pstricks}
\usepackage{epsfig}
\usepackage{multicol}
\usepackage{listings}
\usepackage{enumerate}
\usepackage{hyperref}
\hypersetup{
  bookmarks=true,
  bookmarksopen=true,
  pdfborder=false,
  pdfpagemode=UseNone,
  raiselinks=true,
  pdfhighlight={/P},
  colorlinks,
  citecolor=black,
  filecolor=black,
  linkcolor=black,
  urlcolor=black
}


\usepackage{fontspec}
\usepackage{xunicode}
\usepackage{xltxtra}


% Margins
%\setlength{\textwidth}{24cm}

%\setlength{\voffset}{0in}
%\setlength{\textheight}{6.8in}
% Colors
\definecolor{cbrown}{rgb}{0.49,0.24,0.07}
\definecolor{cmpez}{rgb}{0.92,0.88,0.79}
\definecolor{cmpez2}{rgb}{0.62,0.33,0.34}
\definecolor{cwhite}{rgb}{1,1,1}
\definecolor{cblack}{rgb}{0,0,0}
\definecolor{cred}{rgb}{0.9,0.15,0.15}
\definecolor{clblue}{rgb}{0.57,0.85,0.97}
\definecolor{corange}{rgb}{0.86,0.49,0.18}
\definecolor{clorange}{rgb}{0.95,0.84,0.61}
\definecolor{cyellow}{rgb}{1,0.95,0.16}
\definecolor{cgreen}{rgb}{0.69,0.8,0.31}
\definecolor{clbrown}{rgb}{0.78,0.67,0.43}
\definecolor{cmagenta}{rgb}{0.79,0.05,0.54}
\definecolor{cgray}{rgb}{0.76,0.73,0.63}
\definecolor{bluesite}{rgb}{0.05,0.43,0.67}
\definecolor{redsite}{rgb}{0.90,0.15,0.15}
\definecolor{backsite}{rgb}{0.07,0.07,0.07}
% ----------- from Costas -----------------------------

%---------------   Main Fonts    -----------------%   

% \setmainfont[Mapping=tex-text]{Calibri}
% \setmainfont[Mapping=tex-text]{Times New Roman}
%\setmainfont[Mapping=tex-text]{Droid Serif}
%\setmainfont[Mapping=tex-text]{Cambria}
%\setmainfont[Mapping=tex-text]{cm-unicode}
% \setmainfont[Mapping=tex-text]{Gentium}
% \setmainfont[Mapping=tex-text]{GFS Didot}
% \setmainfont[Mapping=tex-text]{Comic Sans MS}
 \setmainfont[Mapping=tex-text]{Ubuntu}
% \setmainfont[Mapping=tex-text]{Myriad Pro}
% \setmainfont[Mapping=tex-text]{CMU Concrete}
%\setmainfont[Mapping=tex-text]{DejaVu Sans}
%\setmainfont[Mapping=tex-text]{KerkisSans}
%\setmainfont[Mapping=tex-text]{KerkisCaligraphic}
% \setmainfont[Mapping=tex-text]{Segoe Print}
% \setmainfont[Mapping=tex-text]{Gabriola}

%------------------------------------------------%

% ------------ Mathematics Fonts ----------------%
\setmathfont{Asana-Math.ttf}
% -----------------------------------------------%
% Misc
\author{Κασωτάκη Ε. - Σμαραγδάκης Κ.}
\title{www.math24.gr}
% ---------------- Formation --------------------%


\setlength\topmargin{-0.7cm}
\addtolength\topmargin{-\headheight}
\addtolength\topmargin{-\headsep}
\setlength\textheight{26cm}
\setlength\oddsidemargin{-0.54cm}
\setlength\evensidemargin{-0.54cm}
%\setlength\marginparwidth{1.5in}
\setlength\textwidth{17cm}

\RequirePackage[avantgarde]{quotchap}
\renewcommand\chapterheadstartvskip{\vspace*{0\baselineskip}}
\RequirePackage[calcwidth]{titlesec}
\titleformat{\section}[hang]{\bfseries}
{\Large\thesection}{12pt}{\Large}[{\titlerule[0.9pt]}]
%--------------------------------------------------%
\usepackage{draftwatermark}
\SetWatermarkText{www.math24.gr}
\SetWatermarkLightness{0.9}
\SetWatermarkScale{0.6}
%--------------------------------------------------%


\usepackage{xcolor}
\usepackage{amsthm}
\usepackage{framed}
\usepackage{parskip}

\colorlet{shadecolor}{bluesite!20}

\newtheorem*{orismos}{Ορισμός}

\newenvironment{notation}
  {\begin{shaded}\begin{theorem}}
  {\end{theorem}\end{shaded}}




% Document begins
\begin{document}
\pagestyle{fancy}
\fancyhead{}
\fancyfoot{}
\renewcommand{\headrulewidth}{0pt}
\renewcommand{\footrulewidth}{0pt}


\fancyhead[LO,LE]{
 %\textblockcolor{backsite}
 %\begin{textblock}{5}(-2,-0.55)
  %\rule{0cm}{1cm}
 %\end{textblock}
 \textblockcolor{bluesite}
 \begin{textblock}{5}(-1.5,-0.55)
  \rule{0cm}{1cm}
 \end{textblock}
 %\textblockcolor{bluesite}
 %\begin{textblock}{14}(5.5,-0.55)
 % \rule{0cm}{1cm}
 %\end{textblock}
 \begin{textblock}{0}(-1,-0.25)
 \color{cwhite} \begin{Large}www.math24.gr\end{Large}
 \end{textblock}
\begin{textblock}{0}(16,-0.4)
 \color{cwhite} \includegraphics[height=1.5cm]{math24_logo.png}
 \end{textblock}
\textblockcolor{white}
\begin{textblock}{17}(-1,28)
 \color{backsite} \begin{small}Copyright \textcopyright 2013, Κασωτάκη Ε.(ikasotaki@gmail.com) - Σμαραγδάκης Κ.(kesmarag@gmail.com)\end{small}
 \end{textblock}
}
 
\begin{shaded}
\begin{center}
\huge \textbf{Μαθηματικά Α' Γυμνασίου}\\
\end{center} 
\begin{center}
\textbf{Επαναληπτικές ερωτήσεις θεωρίας} \hfill \textbf{}
\end{center}
%\subsection*{}
\end{shaded}
\vspace{2em}

%\begin{itemize}
%\item 1.3: Δυνάμεις φυσικών αριθμών
%\item 1.4: Ευκλείδεια διαίρεση - διαιρετότητα
%\item 1.5: Χαρακτήρες διαιρετότητας - ΜΚΔ - ΕΚΠ - Ανάλυση αριθμού σε γινόμενο πρώτων παραγόντων
%\item Κεφάλαιο 2: Τα κλάσματα
%\item Κεφάλαιο 5: Ποσοστά
%\item 6.3: Ανάλογα ποσά - Ιδιότητες ανάλογων ποσών
%\item 6.5: Προβλήματα αναλογιών
%\item 6.6: Αντιστρόφως ανάλογα ποσά
%\item Β.1.3: Μονάδες μήκους – απόσταση σημείων – μέσο ευθύγραμμου τμήματος 
%\item Β.1.5: Διχοτόμος γωνίας 
%\item Β.1.6: Είδη γωνιών – Κάθετες γωνίες 
%\item Β.1.7: Εφεξής – διαδοχικές γωνίες 
%\item Β.1.8: Παραπληρωματικές και συμπληρωματικές γωνίες – κατακορυφήν 
%\item Β.1.9: Θέσεις ευθειών στο επίπεδο 
%\item Β.1.10: Απόσταση σημείου από ευθεία – Απόσταση παραλλήλων 
%\item Β.2.3: Μεσοκάθετος ευθύγραμμου τμήματος 
%\item Β.2.6: Παράλληλες ευθείες που τέμνονται από μία άλλη ευθεία
%\item Β.3.1: Στοιχεία τριγώνου – Είδη τριγώνων 
%\item Β.3.2: Άθροισμα γωνιών τριγώνου – Ιδιότητες ισοσκελούς τριγώνου 
%\end{itemize}

\section*{Επαναληπτικές Ερωτήσεις Θεωρίας \hfill \small{}}
\begin{enumerate}
\item \textbf{Τι ονομάζεται Ελάχιστο Κοινό Πολλαπλάσιο (ΕΚΠ) δύο ή περισσότερων αριθμών;}\\
      Ελάχιστο Κοινό Πολλαπλάσιο (ΕΚΠ) δύο ή περισσότερων αριθμών ονομάζεται το μικρότερο από τα κοινά 
      πολλαπλάσιά τους που δεν  δεν είναι μηδέν.
\item \textbf{Ποιοι αριθμοί ονομάζονται πρώτοι;}\\
      Ένας αριθμός λέγεται πρώτος όταν έχει διαιρέτες μόνο τον εαυτό του και το $1$.
\item \textbf{Τι ονομάζεται Μέγιστος Κοινός Διαιρέτης (ΜΚΔ) δύο αριθμών;}\\
      Μέγιστος Κοινός Διαιρέτης (ΜΚΔ) δύο αριθμών ονομάζεται ο μεγαλύτερος αριθμός από τους κοινούς 
      τους διαιρέτες.
\item \textbf{Πότε δύο αριθμοί λέγονται πρώτοι μεταξύ τους;}\\
       Δύο αριθμοί λέγονται πρώτοι μεταξύ τους όταν ο μέγιστος κοινός διαιρέτης τους είναι το $1$.
\item \textbf{Πότε ένας φυσικός αριθμός διαιρείται με το $10,100,1000,\cdots$;}\\
      Ένας φυσικός αριθμός διαιρείται με το $10,100,1000,\cdots$ αν λήγει σε ένα, δύο, τρία, $\cdots$ μηδενικά 
      αντίστοιχα.
%-------------------------------------------------------------05
\item \textbf{Πότε ένας φυσικός αριθμός διαιρείται με το $2$;}\\
      Ένας φυσικός αριθμός διαιρείται με το $2$ αν το τελευταίο ψηφίο του είναι $0,2,4,6,8$.
\item \textbf{Πότε ένας φυσικός αριθμός διαιρείται με το $5$;}\\
      Ένας φυσικός αριθμός διαιρείται με το $5$ αν λήγει σε $0$ ή $5$.
\item \textbf{Πότε ένας φυσικός αριθμός διαιρείται με το $3$;}\\
      Ένας φυσικός αριθμός διαιρείται με το $3$ αν το άθροισμα των ψηφίων του διαιρείται με το $3$.
\item \textbf{Πότε ένας φυσικός αριθμός διαιρείται με το $9$;}\\
     Ένας φυσικός αριθμός διαιρείται με το $9$ αν το άθροισμα των ψηφίων του διαιρείται με το $9$.
\item \textbf{Πότε ένας φυσικός αριθμός διαιρείται με το $4$;}\\
      Ένας φυσικός αριθμός διαιρείται με το $4$ αν τα δύο τελευταία ψηφία του σχηματίζουν αριθμό που 
      διαιρείται με το $4$.
%--------------------------------------------------------10
\item \textbf{Πότε ένας φυσικός αριθμός διαιρείται με το $25$;}\\
      Ένας φυσικός αριθμός διαιρείται με το $25$ αν τα δύο τελευταία ψηφία του σχηματίζουν αριθμό που 
      διαιρείται με το $25$.
\item \textbf{Πότε δύο κλάσματα λέγονται ισοδύναμα;}\\
       Δύο κλάσματα λέγονται ισοδύναμα όταν εκφράζουν το ίδιο τμήμα ενός μεγέθους ή ίσων μεγεθών.
\item \textbf{*Τι ονομάζουμε απλοποίηση κλάσματος; }\\
       Απλοποίηση κλάσματος ονομάζεται η διαδικασία με την οποία οι όροι του κλάσματος διαιρούνται με τον 
       ίδιο αριθμό ($\neq 0$) και το ισοδύναμο κλάσμα που προκύπτει έχει μικρότερους όρους από το αρχικό.
\item \textbf{Πότε ένα κλάσμα λέγεται ανάγωγο;}\\
       Ανάγωγο λέγεται το κλάσμα εκείνο που δεν μπορεί να απλοποιηθεί (δηλαδή δεν υπάρχει κοινός διαιρέτης 
       αριθμητή και παρονομαστή).
\item \textbf{Πότε δύο ή περισσότερα κλάσματα λέγονται ομώνυμα;}\\
      Δύο ή περισσότερα κλάσματα λέγονται ομώνυμα όταν έχουν τον ίδιο παρονομαστή.
%------------------------------------------------------15
\item \textbf{Πότε δύο ή περισσότερα κλάσματα λέγονται ετερώνυμα;}\\
       Δύο ή περισσότερα κλάσματα λέγονται ετερώνυμα όταν έχουν διαφορετικούς παρονομαστές.
\item \textbf{Πότε ένα κλάσμα είναι μεγαλύτερο από το $1$;}\\
       Ένα κλάσμα είναι μεγαλύτερο από το $1$ όταν ο αριθμητής είναι μεγαλύτερος από τον παρονομαστή.
\item \textbf{Πότε ένα κλάσμα είναι μικρότερο από το $1$;}\\
       Ένα κλάσμα είναι μικρότερο από το $1$ όταν ο αριθμητής είναι μικρότερος από τον παρονομαστή.
\item \textbf{Πότε ένα κλάσμα ισούται με $1$;}\\
       Ένα κλάσμα ισούται $1$ όταν ο αριθμητής και ο παρονομαστής είναι ίσοι.
\item \textbf{Αν δύο κλάσματα έχουν τον ίδιο παρονομαστή (είναι ομώνυμα) ποιο είναι μεγαλύτερο;}\\
      Αν δύο κλάσματα έχουν τον ίδιο παρονομαστή μεγαλύτερο είναι εκείνο που έχει μεγαλύτερο αριθμητή.
%-------------------------------------------------------20
\item \textbf{Αν δύο κλάσματα έχουν τον ίδιο αριθμητή ποιο είναι μεγαλύτερο;}\\
      Αν δύο κλάσματα έχουν τον ίδιο αριθμητή μεγαλύτερο είναι εκείνο που έχει το μικρότερο παρονομαστή.
\item \textbf{Πώς συγκρίνουμε δύο ετερώνυμα κλάσματα;}\\
      Για να συγκρίνουμε δύο ετερώνυμα κλάσματα τα μετατρέπουμε σε ομώνυμα και συγκρίνουμε τους αριθμητές τους.
\item \textbf{*Πώς προσθέτουμε δύο ή περισσότερα ομώνυμα κλάσματα;}\\
      Προσθέτουμε δύο ή περισσότερα ομώνυμα κλάσματα προσθέτοντας τους αριθμητές τους.
\item \textbf{*Πώς προσθέτουμε ετερώνυμα κλάσματα;}\\
      Για να προσθέσουμε ετερώνυμα κλάσματα πρέπει πρώτα να τα μετατρέψουμε σε ομώνυμα και στη συνέχεια 
      προσθέτουμε τους αριθμητές τους.
\item \textbf{*Πώς αφαιρούμε δύο ομώνυμα κλάσματα;}\\
       Αφαιρούμε δύο ομώνυμα κλάσματα αφαιρώντας τους αριθμητές τους.
%--------------------------------------------------------25
\item \textbf{*Πώς αφαιρούμε δύο ετερώνυμα κλάσματα;}
       Αφαιρούμε δύο ομώνυμα κλάσματα αφού πρώτα τα μετατρέψουμε σε ομώνυμα και τότε στη συνέχεια 
       αφαιρούμε τους αριθμητές τους.
\item \textbf{*Πώς υπολογίζουμε το γινόμενο δύο κλασμάτων;}\\
       Το γινόμενο δύο κλασμάτων είναι το κλάσμα που έχει αριθμητή το γινόμενο των αριθμητών και 
       παρονομαστή το γινόμενο των παρονομαστών.
\item \textbf{*Πώς υπολογίζουμε το γινόμενο ενός φυσικού αριθμού επί ένα κλάσμα;}\\
      Το γινόμενο ενός φυσικού αριθμού επί ένα κλάσμα είναι το κλάσμα με αριθμητή το γινόμενο του αριθμητή 
      επί το φυσικό αριθμό και με τον ίδιο παρονομαστή.
\item \textbf{Πότε δύο κλάσματα λέγονται αντίστροφα;}\\
      Δύο κλάσματα λέγονται αντίστροφα όταν έχουν γινόμενο $1$.
\item \textbf{Πότε δύο αριθμοί λέγονται αντίστροφοι;}\\
      Δύο αριθμοί λέγονται αντίστροφοι όταν το γινόμενό τους ισούται με $1$.
%-------------------------------------------------------30
\item \textbf{*Πώς υπολογίζουμε το πηλίκο δύο κλασμάτων;}\\
      Για να διαιρέσουμε δύο κλάσματα πολλαπλασιάζουμε το διαιρετέο με το αντίστροφο του διαιρέτη.
\item \textbf{Τι ονομάζονται σύνθετο κλάσμα;}\\
     Σύνθετο κλάσμα ονομάζεται το κλάσμα του οποίου ένας τουλάχιστον όρος του είναι κλάσμα.
\item \textbf{*Πώς ονομάζεται και με τι ισούται το σύμβολο $α\%$;}\\
      Το σύμβολο $α\%$ ονομάζεται ποσοστό επί τοις εκατό ή απλούστερα ποσοστό και είναι ίσο με $\dfrac{α}{100}$.
\item \textbf{*Πώς υπολογίζουμε το ποσοστό $α\%$ ενός αριθμού $β$;}\\
      Υπολογίζουμε το ποσοστό $α\%$ ενός αριθμού $β$ με το γινόμενο $\dfrac{α}{100} \cdot β$.
\item \textbf{Πότε δύο ποσά λέγονται ανάλογα;}\\
       Δύο ποσά λέγονται ανάλογα εάν μεταβάλλονται με τέτοιο τρόπο που όταν οι τιμές του ενός πολλαπλασιάζονται 
       με έναν αριθμό, τότε οι αντίστοιχες τιμές του άλλου πολλαπλασιάζονται με τον ίδιο αριθμό.
%------------------------------------------------------35
\item \textbf{Τι ονομάζουμε συντελεστή αναλογίας δύο ανάλογων ποσών;}\\
      Συντελεστής αναλογίας δύο ανάλογων ποσών ονομάζεται το σταθερό πηλίκο τους.
\item \textbf{Πότε δύο μεγέθη ονομάζονται αντιστρόφως ανάλογα;}\\
       Δύο μεγέθη ονομάζονται αντιστρόφως ανάλογα στην περίπτωση που η μεταβολή τους είναι τέτοια ώστε:
       όταν το ένα μέγεθος πολλαπλασιάζεται επί έναν αριθμό, το άλλο διαιρείται με τον ίδιο αριθμό.
\item \textbf{Ποια είναι η μονάδα μέτρησης του μήκους και ποιες είναι οι υποδιαιρέσεις και ποια τα πολλαπλάσιά της;}\\
      Η μονάδα μέτρησης του μήκους είναι το $1m$ (μέτρο)\\
      πολλαπλάσια: $1km=1000m$\\
      υποπολλαπλάσια: $1dm=\dfrac{1}{10},1cm=\dfrac{1}{100},1mm=\dfrac{1}{1000}$
\item \textbf{*Τι ονομάζουμε απόσταση δύο σημείων;}\\
      Απόσταση δύο σημείων $Α$ και $Β$ λέγεται το μήκος του ευθύγραμμου τμήματος $ΑΒ$ που τα ενώνει.
\item \textbf{Τι ονομάζουμε μέσο ενός ευθύγραμμου τμήματος;}\\
      Μέσο ενός ευθύγραμμου τμήματος $ΑΒ$ ονομάζουμε το σημείο $Μ$ του τμήματος που ισαπέχει από τα άκρα του.
%------------------------------------------------------40
\item \textbf{*Ποια είναι η μονάδα μέτρησης της γωνίας;}\\
      Μονάδα μέτρησης της γωνίας είναι η μοίρα.\\
      υποπολλαπλάσια: $1$\textdegree$=60'$ και $1'=60''$
\item \textbf{Πότε δύο γωνίες είναι ίσες;}\\
      Δύο γωνίες είναι ίσες όταν έχουν το ίδιο μέτρο.
\item \textbf{Τι ονομάζουμε διχοτόμος μιας γωνίας;}\\
      Διχοτόμος μιας γωνίας ονομάζεται η ημιευθεία που έχει αρχή την κορυφή της γωνίας και τη χωρίζει σε δύο ίσες γωνίες.
\item \textbf{Ποια γωνία ονομάζεται ορθή;}\\
      Ορθή γωνία λέγεται η γωνία της οποίας το μέτρο είναι ίσο με $90$\textdegree.
\item \textbf{Ποια γωνία ονομάζεται οξεία;}\\
      Οξεία γωνία λέγεται κάθε γωνία που έχει μέτρο μικρότερο από $90$\textdegree.
%--------------------------------------------------------------45
\item \textbf{Ποια γωνία ονομάζεται αμβλεία;}\\
      Αμβλεία γωνία λέγεται κάθε γωνία που έχει μέτρο μεγαλύτερο των $90$\textdegree και μικρότερο των $180$\textdegree.
\item \textbf{Ποια γωνία ονομάζεται ευθεία;}\\
      Μια γωνία λέγεται ευθεία όταν το μέτρο της ισούται με $180$\textdegree.
\item \textbf{Ποια γωνία ονομάζεται κυρτή;}\\
      Μια γωνία ονομάζεται κυρτή όταν έχει μέτρο μεγαλύτερο των $180$\textdegree και μικρότερο των $360$\textdegree.
\item \textbf{Ποια γωνία ονομάζεται μηδενική;}\\
      Μηδενική γωνία λέγεται η γωνία της οποίας το μέτρο είναι ίσο με $0$\textdegree.
\item \textbf{Ποια γωνία ονομάζεται πλήρης;}\\
      Μια γωνία λέγεται πλήρης όταν το μέτρο της είναι ίσο με $360$\textdegree.
%-------------------------------------------------------------50
\item \textbf{Πότε δύο ευθείες είναι κάθετες;}\\
      Δύο ευθείες είναι κάθετες όταν οι γωνίες που σχηματίζουν αυτές τεμνόμενες είναι ορθές.
\item \textbf{*Πότε δύο ευθύγραμμα τμήματα λέμε ότι είναι κάθετα;}\\
      Δύο ευθύγραμμα τμήματα λέμε ότι είναι κάθετα όταν βρίσκονται σε δύο κάθετες ευθείες.
\item \textbf{*Πότε δύο ημιευθείες λέμε ότι είναι κάθετες;}\\
      Δύο ημιευθείες λέμε ότι είναι κάθετες όταν βρίσκονται σε δύο κάθετες ευθείες.
\item \textbf{Πότε δύο γωνίες ονομάζονται εφεξής;}\\
       Εφεξής ονομάζονται δύο γωνίες που έχουν την ίδια κορυφή, μία κοινή πλευρά και δεν έχουν κανένα 
       άλλο κοινό σημείο.
\item \textbf{Πότε δύο γωνίες ονομάζονται διαδοχικές;}\\
      Διαδοχικές γωνίες λέγονται περισσότερες από δύο γωνίες που, βρίσκονται στο ίδιο επίπεδο και, καθεμία από 
      αυτές είναι εφεξής γωνία με την προηγούμενη ή την επόμενή της.
%-------------------------------------------------------------55
\item \textbf{Πότε δύο γωνίες ονομάζονται παραπληρωματικές;}
      Παραπληρωματικές γωνίες ονομάζονται δύο γωνίες που έχουν άθροισμα $180$\textdegree.
\item \textbf{Πότε δύο γωνίες ονομάζονται συμπληρωματικές;}\\
      Συμπληρωματικές γωνίες ονομάζονται δύο γωνίες που έχουν άθροισμα $90$\textdegree
\item \textbf{Ποιες γωνίες ονομάζονται κατακορυφήν;}\\
      Κατακορυφήν γωνίες ονομάζονται δύο γωνίες που έχουν την κορυφή τους κοινή και τις πλευρές τους 
      αντικείμενες ημιευθείες.
\item \text{Πότε δύο ευθείες είναι παράλληλες;}\\
       Δύο ευθείες του ίδιου επιπέδου λέγονται παράλληλες αν δεν έχουν κοινό σημείο όσο κι αν προεκταθούν.
\item \textbf{Πόε δύο ευθείες ονομάζονται τεμνόμενες;}\\
      Δύο ευθείες του ίδιου επιπέδου που έχουν ένα κοινό σημείο ονομάζονται τεμνόμενες (και το κοινό τους 
      σημείο ονομάζεται σημείο τομής των δύο ευθειών).
%---------------------------------------------------------------60
\item \textbf{*Πότε δύο ευθύγραμμα τμήματα λέγονται παράλληλα;}\\
      Δύο ευθύγραμμα τμήματα λέγονται παράλληλα όταν βρίσκονται σε παράλληλες ευθείες.
\item \textbf{*Τι ονομάζουμε απόσταση ενός σημείου $Α$ από μία ευθεία $ε$;}\\
      Απόσταση του σημείου $Α$ από μία ευθεία $ε$ ονομάζεται το μήκος του κάθετου ευθυγράμμου τμήματος $ΑΑ_{0}$
      από το σημείο $Α$ προς την ευθεία $ε$.
\item \textbf{Τι ονομάζουμε απόσταση δύο παραλλήλων ευθειών;}\\
      Απόσταση δύο παραλλήλων ευθειών λέγεται το μήκος οποιουδήποτε ευθυγράμμου τμήματος που είναι κάθετο στις δύο ευθείες 
      και έχει τα άκρα του σ' αυτές.
\item \textbf{Τι ονομάζουμε μεσοκάθετο ευθυγράμμου τμήματος;}\\
      Μεσοκάθετος ευθυγράμμου τμήματος λέγεται η ευθεία που είναι κάθετη προς αυτό και διέρχεται από το μέσον του.
\item \textbf{*Για τις παράλληλες ευθείες $ε_{1}$ και $ε_{2}$ που τέμνονται από μία τρίτη ευθεία $δ$,
       ποιες γωνίες ονομάζονται "εντός", ποιες "εκτός", ποιες "επί τα αυτά" και ποιες "εναλλάξ";}\\
       \begin{itemize}
        \item \textbf{"εντός"} ονομάζονται οι γωνίες που βρίσκονται ανάμεσα στις $ε_{1}$ και $ε_{2}$ και όλες 
              οι άλλες \textbf{"εκτός"}
        \item \textbf{"επί τα αυτά"} ονομάζονται οι γωνίες που βρίσκονται προς το ίδιο μέρος της ευθείας $δ$
        \item \textbf{"εναλλάξ"} ονομάζονται δύο γωνίες που βρίσκονται η μία στο ένα κι η άλλη στο άλλο ημιεπίπεδο 
               της ευθείας $δ$
       \end{itemize}

%-----------------------------------------------------------------65
\item \textbf{Ποια είναι τα κύρια στοιχεία ενός τριγώνου;}\\
      Τα κύρια στοιχεία ενός τριγώνου είναι: 
      \begin{itemize}
       \item οι τρεις κορυφές του
       \item οι τρεις γωνίες του
       \item οι τρεις πλευρές του
      \end{itemize}

\item \textbf{Πότε ένα τρίγωνο ονομάζεται ορθογώνιο;}\\
      Ένα τρίγωνο ονομάζεται ορθογώνιο όταν έχει μία γωνία ορθή.
\item \textbf{Πότε ένα τρίγωνο ονομάζεται αμβλυγώνιο;}\\
       Ένα τρίγωνο ονομάζεται αμβλυγώνιο όταν έχει μία γωνία αμβλεία.
\item \textbf{Πότε ένα τρίγωνο ονομάζεται οξυγώνιο;}\\
       Ένα τρίγωνο ονομάζεται οξυγώνιο όταν και οι τρεις γωνίες του είναι οξείες.
\item \textbf{Πότε ένα τρίγωνο ονομάζεται ισόπλευρο;}\\
       Ένα τρίγωνο ονομάζεται ισόπλευρο όταν και οι τρεις πλευρές του είναι ίσες.
%---------------------------------------------------------------70
\item \textbf{Πότε ένα τρίγωνο ονομάζεται ισοσκελές;}\\
       Ένα τρίγωνο ονομάζεται ισοσκελές όταν έχει δύο πλευρές ίσες.
\item \textbf{Πότε ένα τρίγωνο ονομάζεται σκαληνό;}\\
       Ένα τρίγωνο ονομάζεται σκαληνό όταν όλες οι πλευρές του είναι άνισες.
\item \textbf{*Ποια είναι τα δευτερεύοντα στοιχεία ενός τριγώνου;}\\
      Τα δευτερεύοντα στοιχεία ενός τριγώνου είναι:
      \begin{itemize}
       \item Διάμεσος 
       \item Ύψος
       \item Διχοτόμος
      \end{itemize}

\item \textbf{Τι ονομάζουμε διάμεσο ενός τριγώνου;}\\
      Διάμεσος λέγεται το ευθύγραμμο τμήμα που ενώνει την κορυφή ενός τριγώνου με το μέσο της απέναντι πλευράς.
\item \textbf{Τι ονομάζουμε ύψος τριγώνου;}\\
      Ύψος ονομάζουμε το ευθύγραμμο τμήμα που φέρνουμε από την κορυφή ενός τριγώνου κάθετο στην ευθεία της 
      απέναντι πλευράς.
%---------------------------------------------------------------75
\item \textbf{Τι ονομάζουμε διχοτόμο του τριγώνου;}\\
      Διχοτόμο του τριγώνου ονομάζουμε το ευθύγραμμο τμήμα της διχοτόμου  μιας γωνίας ενός τριγώνου που φέρνουμε 
      από μια κορυφή και καταλήγει στην απέναντι πλευρά.
\item \textbf{Με τι ισούται το άθροισμα των γωνιών κάθε τριγώνου;}\\
      Το άθροισμα των γωνιών κάθε τριγώνου είναι ίσο με $180$\textdegree.
\end{enumerate}



\end{document}
